\section{\heiti 结论}

本文选择上下文规则优先的混合分类策略，其根本依据在于：能源财务模型的公式业务含义由行列标题决定而非由公式语法决定——消融实验中同一SUM函数在不同行代表完全不同的业务类型，这一结构性特征使得纯语法方法从根本上不适用于本场景，而非仅仅在准确率上略有不足。

消融实验揭示的核心规律是：跨业务Sheet现象（单一Sheet内嵌多个业务模块）构成了规则方法的天然边界。本模型时间序列Sheet内嵌折旧时间轴逻辑，导致28/35条错误集中于此。这一发现说明，知识图谱构建中的语义歧义问题往往来自模型设计层面的约定，而非方法本身的局限——这对财务模型的规范化设计具有直接的实践指导意义。

知识图谱的工程价值不仅体现在查询响应时间（IRR路径追溯<0.5秒）上，更体现在使财务模型从单点专家知识转变为可查询、可审计的组织知识资产——这是传统Excel文件无法实现的根本转变。当前模型包含289个含错误公式，这些错误在人工审计中极难全面检出，而知识图谱的依赖图可以直接定位错误节点对下游指标的影响范围，这一能力在本文实验中尚未充分展开，是后续研究的重点方向之一。

本研究的局限在于两点：其一，研究对象为单一抽水蓄能项目模型，风电、光伏、核电等类型的财务模型是否存在相同的跨业务Sheet现象尚待验证；其二，LLM增强分类模块当前仅有框架设计，21.7\%的高歧义公式的实际分类效果有待实验验证。后续工作将重点收集多类型能源财务模型，构建多案例对比基准，同时探索基于LLM的跨业务公式语义消歧方法。

\vspace {3mm}
\zihao{5}{
	\noindent \textsf{致\quad 谢}\quad \textit{感谢提供真实抽水蓄能电站财务评价模型数据的研究合作方，以及在数据标注和实验验证过程中给予支持的团队成员。}}

\vspace {5mm}
\centerline
{\zihao{5}\textsf{参~考~文~献}}
\zihao{5-} \addtolength{\itemsep}{-1em}
\vspace {1.5mm}
\bibliographystyle{gbt7714-numerical}
\bibliography{ref.bib}
